% Strict but not uniform: the spectral gap at every finite extent.
% Lane: O (operator bridge), modules YangMills/OS/**.
% REGIME: tricotomy; NO scores, NO desk language, NO commit chronology beyond
% the reproducibility freeze.
\documentclass[11pt]{article}
\usepackage[margin=1.1in]{geometry}
\usepackage{amsmath,amssymb,amsthm}
\usepackage{booktabs,longtable,array}
\usepackage[colorlinks=true,linkcolor=blue,urlcolor=blue,citecolor=blue]{hyperref}

\newtheorem{theorem}{Theorem}[section]
\newtheorem{lemma}[theorem]{Lemma}
\newtheorem{proposition}[theorem]{Proposition}
\newtheorem{corollary}[theorem]{Corollary}
\newtheorem{remark}[theorem]{Remark}
\theoremstyle{definition}
\newtheorem{definition}[theorem]{Definition}

\emergencystretch=3em
\tolerance=2000

\newcommand{\anchor}{ff24cf46cfc8859dffa8b2f18252763e0c75def9}
\newcommand{\osline}[3]{%
  \href{https://github.com/lluiseriksson/THE-ERIKSSON-PROGRAMME/blob/\anchor/YangMills/OS/#1.lean\#L#2}{\nolinkurl{#3}}}
\newcommand{\repofile}[2]{%
  \href{https://github.com/lluiseriksson/THE-ERIKSSON-PROGRAMME/blob/\anchor/#1}{\nolinkurl{#2}}}
\newcommand{\R}{\mathbb{R}}
\newcommand{\Z}{\mathbb{Z}}
\newcommand{\SU}{\mathrm{SU}}
\newcommand{\Om}{\Omega}

\title{Strict but Not Uniform:\\
a Machine-Checked Spectral Gap at Every Finite Extent\\
of the Coupled Slice\\
\large peripheral separation without an equality case,\\
and the modulus the theorem does not provide}
\author{Lluis Eriksson}
\date{}

\begin{document}
\maketitle

\begin{abstract}
A companion paper supplied the vacuum of the coupled $\Z_2$ slice at every
spatial extent: a strictly positive eigenvector, unique up to scale, carrying
the spectral radius \cite{ErikssonPerron}.  It listed \emph{peripheral
separation} as out of scope, and that omission is not cosmetic --- without it
$|\mu|\le\lambda$ leaves $\mu=-\lambda$ open, and no gap follows at all.

This paper closes it, and then draws the consequence that matters, which is
negative.

\textbf{Proved.}  For a strictly positive kernel on a finite nonempty type,
$-\lambda$ is not an eigenvalue, hence every real eigenvalue other than the
Perron eigenvalue is \emph{strictly} smaller in absolute value.  Specialised to
the coupled slice: at every extent $L$, every $\beta$, and every strictly
positive source weight, the transfer operator has a strict spectral gap.  The
proof of peripheral separation avoids the equality case of the triangle
inequality, which is where the classical argument spends its effort: writing
$u=|w|$, $p=u-w$ and $q=u+w$, one gets $Ap=\lambda q$ and $Aq=\lambda p$, so a
nonzero $p$ would make $Ap$ strictly positive, hence $q$ strictly positive,
hence $w$ nonnegative, hence $p=0$.

The separation is then extended from the real eigenvalues to \emph{all} of them.
The coupled kernel is conjugate by a positive diagonal to its symmetrised form,
which is symmetric; and a real symmetric kernel has real eigenvalues, by a
computation that pairs the eigenvector against its image twice and is two
rearrangements of a double sum.  So there are no complex peripheral eigenvalues
left to exclude, and the strict gap is a statement about the whole spectrum.  We
also deliver the vacuum in Euclidean normalisation, $\|\Om\|=1$ with
$T\Om=\Om$.

\textbf{Not proved, and this is the title.}  The gap is \emph{strict}, not
\emph{quantitative}: the theorem provides no modulus of separation, and in
particular nothing uniform in $L$.  That is not a limitation we expect to remove
by working harder in this direction.  Direct numerical diagonalisation shows the
subdominant ratio running $0.9205,\,0.9829,\,0.9964,\,0.9992$ at $L=2,3,4,5$ for
$\beta=0.8$, $\gamma=1.2$ --- collapsing towards $1$.  That computation is
reported as measured and unproved, and no theorem here depends on it; its role
is to say that a paper reporting only the positive half would be reporting the
half that does not matter.

Nothing in this paper is a claim about $\SU(N)$, the continuum limit, or the
Yang--Mills mass gap.
\end{abstract}

\section{Introduction}

\subsection{What was missing, and why it was not cosmetic}

\cite{ErikssonPerron} proved that the coupled slice has a vacuum at every finite
extent --- strictly positive, unique up to scale, carrying the spectral radius
--- and listed among its exclusions that \emph{peripheral separation} was not
proved.  Read quickly, that looks like a technical refinement.  It is not.  The
domination statement is $|\mu|\le\lambda$; the case $\mu=-\lambda$ saturates it.
Until that case is excluded, the second-largest modulus may equal the largest,
and \emph{no gap statement exists at all}, at any extent.

So the first half of this paper is not an improvement of the previous one.  It
is the step without which the previous one supports no spectral conclusion.

\subsection{What is proved}

Throughout $A$ is a strictly positive kernel on a finite nonempty type, $v>0$
with $Av=\lambda v$.

\begin{itemize}
\item $-\lambda$ is not an eigenvalue of $A$;
\item hence for every real eigenpair $(\mu,w)$ with $w\neq0$ and $\mu\neq\lambda$,
      $\;|\mu|<\lambda$ --- a strict gap;
\item specialised: at every extent $L$, every $\beta$, and every strictly
      positive source weight, the coupled transfer operator has a strict gap;
\item the vacuum in Euclidean normalisation: $\|\Om\|=1$ and $T\Om=\Om$;
\item the symmetrised kernel is symmetric, and strictly positive;
\item \emph{every} eigenvalue of the coupled kernel is real --- the kernel is
      conjugate by a positive diagonal to the symmetrised one --- so the strict
      separation above is a statement about the whole spectrum, not only about
      its real part.
\end{itemize}

\subsection{What is deliberately not claimed}

\begin{enumerate}
\item \textbf{No modulus of separation, and nothing uniform in $L$.}  The gap
theorem is a strict inequality with no quantitative content.  This is the
paper's own headline and Section~\ref{sec:measured} says what the numbers show.
\item \textbf{No proof that a volume-uniform gap fails.}  The collapse in
Section~\ref{sec:measured} is measured on a finite grid, is labelled as such,
and carries no theorem.  What \emph{is} established is that this paper does not
provide a uniform bound --- not that none exists.
\item \textbf{Algebraic simplicity} is not proved, and neither is
Perron--Frobenius for merely irreducible kernels.
\item \textbf{Nothing about $\SU(N)$, the continuum limit, or the Clay problem}
\cite{Clay}.
\end{enumerate}

\section{The positive left eigenvector, and what it buys}

\begin{lemma}[Left eigenvector]\label{lem:left}
$A^{\mathsf T}$ is strictly positive, so it too has a strictly positive
eigenvector $\varphi$, and its eigenvalue equals $\lambda$.
\end{lemma}

\osline{PerronGap}{74}{exists\_pos\_left\_eigenvector},
\osline{PerronGap}{82}{left\_eigenvalue\_eq}.  Equality of the two
eigenvalues follows from pairing: $\langle\varphi,Av\rangle$ evaluates to
$\lambda\langle\varphi,v\rangle$ one way and $\nu\langle\varphi,v\rangle$ the
other, and $\langle\varphi,v\rangle>0$.

\begin{lemma}[A nonnegative sub-eigenvector is an eigenvector]\label{lem:sub}
If $u\ge0$ and $Au\ge\lambda u$ then $Au=\lambda u$.
\end{lemma}

\begin{proof}
Let $z=Au-\lambda u\ge0$.  Then $\langle\varphi,z\rangle
=\langle A^{\mathsf T}\varphi,u\rangle-\lambda\langle\varphi,u\rangle=0$, and a
nonnegative vector with zero positive-weighted sum vanishes.
\end{proof}

\osline{PerronGap}{114}{subeigen\_eq}.  This is the standard device, and it
is what makes the next section cheap.

\section{Peripheral separation, without an equality case}

\begin{theorem}[$-\lambda$ is not an eigenvalue]\label{thm:periph}
There is no $w\neq0$ with $Aw=-\lambda w$.
\end{theorem}

\begin{proof}
Let $u=|w|$ entrywise.  From $\lambda u_i=|\,(Aw)_i\,|\le(Au)_i$ we get
$Au\ge\lambda u$, so $Au=\lambda u$ by Lemma~\ref{lem:sub}.  Put $p=u-w\ge0$ and
$q=u+w\ge0$.  Then
\[
  Ap = Au - Aw = \lambda u + \lambda w = \lambda q,
  \qquad
  Aq = Au + Aw = \lambda u - \lambda w = \lambda p .
\]
Suppose $p\neq0$.  Then $Ap>0$ entrywise, so $\lambda q>0$, so $q_i>0$ for every
$i$; but $q_i=|w_i|+w_i$, which vanishes wherever $w_i<0$, so $w\ge0$ --- and
then $u=w$, i.e.\ $p=0$, a contradiction.  Hence $p=0$, so $w=u\ge0$, so
$Aw=\lambda w$; with $Aw=-\lambda w$ and $\lambda>0$ this forces $w=0$.
\end{proof}

\osline{PerronGap}{163}{neg\_not\_eigenvalue}.

\begin{remark}[Why this route was taken]
The classical proof of peripheral separation goes through the equality case of
the triangle inequality: from $\bigl|\sum_j A_{ij}w_j\bigr|=\sum_j
A_{ij}|w_j|$ one concludes that the $w_j$ share a phase.  Formalising an
equality case is markedly more work than formalising an inequality, and it is
avoidable here.  The pair $(p,q)$ carries the same information in a form that
only ever needs \emph{strict positivity of the image of a nonnegative nonzero
vector} --- a lemma the companion development already had.
\end{remark}

\section{The gap}

\begin{theorem}[Strict gap]\label{thm:gap}
If $Aw=\mu w$ with $w\neq0$ real and $\mu\neq\lambda$, then $|\mu|<\lambda$.
\end{theorem}

\begin{proof}
$|\mu|\le\lambda$ by the domination theorem of \cite{ErikssonBirkhoff}.  If
$|\mu|=\lambda$ then $\mu=\lambda$, excluded by hypothesis, or $\mu=-\lambda$,
excluded by Theorem~\ref{thm:periph}.
\end{proof}

\osline{PerronGap}{246}{abs\_lt\_of\_ne\_perron}, and for the coupled slice
at every extent \osline{PerronGap}{414}{coupled\_gap\_of\_sourceWeight}.

\begin{remark}[What the statement does and does not contain]
Theorem~\ref{thm:gap} is a strict inequality between two real numbers.  It
contains no $\varepsilon$, no dependence on $L$, and therefore no information
whatever about how the separation behaves as the extent grows.  Every use of it
that sounds quantitative is a misuse.
\end{remark}

\section{The two interfaces the next module needs}

\begin{theorem}[Euclidean-normalised vacuum]\label{thm:unit}
With $T=\lambda^{-1}A$ there is $\Om>0$ with $\sum_i\Om_i^2=1$ and $T\Om=\Om$.
\end{theorem}

\osline{PerronGap}{270}{eucNorm},
\osline{PerronGap}{279}{unitVacuum},
\osline{PerronGap}{286}{unitVacuum\_norm},
\osline{PerronGap}{303}{unitVacuum\_fixed}.  The companion paper normalised
on the simplex, which is what its existence proof produced; the
Osterwalder--Seiler side asks for unit Euclidean norm, and the passage between
them is a rescaling.

\begin{theorem}[The symmetrised kernel is symmetric]\label{thm:sym}
$\widetilde K_w(\sigma,\tau)=w(\sigma)^{1/2}K(\sigma,\tau)w(\tau)^{1/2}$
satisfies $\widetilde K_w(\sigma,\tau)=\widetilde K_w(\tau,\sigma)$, and is
strictly positive when $w$ is.
\end{theorem}

\osline{PerronGap}{328}{symWeighted\_symm},
\osline{PerronGap}{335}{symWeighted\_pos}.  It is symmetric because the
decoupled kernel is.

\section{Every eigenvalue is real, so the gap is about the whole spectrum}
\label{sec:real}

Theorem~\ref{thm:gap} separates the \emph{real} eigenvalues.  For a real
non-symmetric kernel that is not yet a statement about the spectrum: nothing so
far excludes a peripheral pair $\mu=\lambda e^{\pm i\theta}$.  The symmetry just
established closes exactly that gap in the argument.

\begin{theorem}[A real symmetric kernel has real eigenvalues]\label{thm:symreal}
If $S$ is real and symmetric and $Sy=\mu y$ with $y\neq0$ complex, then
$\bar\mu=\mu$.
\end{theorem}

\begin{proof}
Pair $y$ against $Sy$: the result is $\mu N$ with $N=\sum_i|y_i|^2>0$.
Conjugating and using that $S$ is real and symmetric returns the same quantity,
so $\bar\mu N=\mu N$.
\end{proof}

\osline{PerronGap}{349}{eigenvalue\_real\_of\_symm}.  No spectral theory
is invoked; the computation is two rearrangements of a double sum.

\begin{theorem}[Every eigenvalue of the coupled kernel is real]\label{thm:allreal}
$K_w$ is conjugate to $\widetilde K_w$ by the positive diagonal
$\mathrm{diag}(w^{1/2})$, so every complex eigenvalue of $K_w$ is an eigenvalue
of $\widetilde K_w$, hence real.
\end{theorem}

\osline{PerronGap}{428}{sourceWeighted\_eq\_sym},
\osline{PerronGap}{443}{coupled\_eigenvalue\_real}.

\begin{corollary}\label{cor:whole}
For the coupled slice, Theorem~\ref{thm:gap} separates \emph{every} eigenvalue
from $\lambda$, not merely the real ones.
\end{corollary}

\begin{remark}[What this does not turn into]
Making the separation a statement about the whole spectrum changes nothing about
its \emph{quantitative} content, which remains empty.  Self-adjointness on an
inner-product space, and the variational characterisation that would supply a
modulus, still belong to the module that consumes this one.
\end{remark}

\section{What is measured and not proved}\label{sec:measured}

Direct dense numerical diagonalisation of the $2^L\times2^L$ coupled kernel
gives, for $\beta=0.8$ and $\gamma=1.2$, subdominant ratios
\[
  0.9205,\quad 0.9829,\quad 0.9964,\quad 0.9992
  \qquad\text{at } L=2,3,4,5,
\]
and for $\beta=0.3$, $\gamma=0.4$,
\[
  0.4665,\quad 0.5518,\quad 0.5938,\quad 0.6159 .
\]
The first family collapses towards $1$; the second does not, and the parameters
of the second satisfy $\sinh 2\beta\,\sinh 2\gamma<1$.  All of this is
\textbf{verified} in the numerical sense and \textbf{not proved}; the
identification of that threshold with the Kramers--Wannier line of the
two-dimensional Ising transfer matrix is standard literature and is not
established here.  The probe and its transcript are in the repository.

No theorem of this paper depends on this section.  Its function is to prevent a
misreading: Theorem~\ref{thm:gap} holds at every finite extent, and the numbers
show that this is compatible with the separation vanishing in the limit.  A
strict gap at each $L$ and a uniform gap in $L$ are different statements, and
only the first is proved here.

\section{Formalisation notes}

\paragraph{An inequality instead of an equality case.}  The one design decision
that shaped the module is described in the remark after
Theorem~\ref{thm:periph}: the $(p,q)$ pair replaces the equality case of the
triangle inequality by two applications of a positivity lemma that already
existed.  The saving is not marginal --- an equality case would have required
either complex-phase reasoning or a sign decomposition of the index set.

\paragraph{The left eigenvector was free.}  Existence for $A^{\mathsf T}$ is the
companion theorem applied to a transposed kernel, and equality of the two
eigenvalues is one pairing.  That single object converts the sub-eigenvector
inequality into an equality, which is the hinge of the whole argument.

\paragraph{Normalisation conventions are interfaces, not mathematics.}  The
companion paper produced $\sum\Om=1$ because its existence proof lives on the
simplex; this one adds $\|\Om\|=1$ because that is what the consuming side
asks.  Both are one rescaling from each other, and neither is a theorem about
the model.  They are recorded as separate endpoints so that no later module has
to guess which convention is in force.

\section{Reproducibility}\label{sec:repro}

All artifacts are at source checkpoint \texttt{\anchor} on branch \texttt{main}
of \url{https://github.com/lluiseriksson/THE-ERIKSSON-PROGRAMME}.  Every link
was resolved against that commit, and checked to point at an actual declaration
line, before compilation.

\begin{center}
\begin{tabular}{ll}
\toprule
Quantity & Value \\
\midrule
Full core build & 8428 jobs, success \\
Oracle commands & 2582 \\
\quad with axiom dependencies & 2559 \\
\quad axiom-free & 23 \\
Distinct axioms across all outputs & \texttt{propext}, \texttt{Quot.sound}, \texttt{Classical.choice} \\
Occurrences of \texttt{sorryAx} & 0 \\
Project axioms & 0 \\
Errors & 0 \\
\bottomrule
\end{tabular}
\end{center}

The numerical probe is
\repofile{scripts/probe\_spatial\_birkhoff.py}{scripts/probe\_spatial\_birkhoff.py}
with transcript
\repofile{docs/O-LANE-PROBE-TRANSCRIPT-20260728.txt}{docs/O-LANE-PROBE-TRANSCRIPT-20260728.txt}.
Verdicts are in
\repofile{docs/VERIFICATION-LEDGER.md}{docs/VERIFICATION-LEDGER.md}.

\begin{thebibliography}{9}

\bibitem{ErikssonPerron}
L.~Eriksson,
\emph{The Vacuum Was Never Absent: A Machine-Checked Perron Theorem for Strictly
Positive Kernels, and the Coupled-Slice Vacuum at Every Spatial Extent},
viXra, 2026.

\bibitem{ErikssonBirkhoff}
L.~Eriksson,
\emph{Blind to the Coupling: a Second Machine-Checked Obstruction at Spatial
Extent}, viXra, 2026.

\bibitem{ErikssonSpatial}
L.~Eriksson,
\emph{Where the Elementary Reconstruction Stops: Spatial Coupling Breaks the
Uniform Vacuum, Machine-Checked}, viXra, 2026.

\bibitem{Seneta}
E.~Seneta,
\emph{Non-negative Matrices and Markov Chains},
2nd ed., Springer, 2006.

\bibitem{Lean4}
L.~de~Moura and S.~Ullrich,
\emph{The Lean 4 theorem prover and programming language},
CADE-28, Lecture Notes in Computer Science \textbf{12699} (2021), 625--635.

\bibitem{Mathlib}
The mathlib Community,
\emph{The Lean mathematical library},
CPP 2020, 367--381.

\bibitem{Clay}
A.~Jaffe and E.~Witten,
\emph{Quantum Yang--Mills theory},
Clay Mathematics Institute Millennium Prize Problem description, 2000.

\end{thebibliography}

\end{document}
